{
\renewcommand{\arraystretch}{1.5}
\setlist[enumerate]{leftmargin=*, labelsep=0.5em}

\begin{tabularx}{\textwidth}{|>{\raggedright\arraybackslash}X|>{\raggedright\arraybackslash}X|}
\hline
\textbf{Tên Use case} & Xóa bộ thẻ ghi nhớ \\
\hline
\textbf{ID} & UC-0034 \\
\hline
\textbf{Tác nhân} & Người dùng đã xác thực (chủ sở hữu bộ thẻ) \\
\hline
\textbf{Mô tả} & Người dùng xóa bộ thẻ ghi nhớ cùng tất cả các thẻ bên trong \\
\hline
\textbf{Điều kiện tiên quyết} & Bộ thẻ tồn tại, người dùng là tác giả \\
\hline
\textbf{Điều kiện sau} & Bộ thẻ và tất cả các thẻ ghi nhớ bên trong bị xóa (cascade delete) \\
\hline
\textbf{Kích hoạt} & Người dùng chọn xóa bộ thẻ \\
\hline
\textbf{Luồng chính} &
\begin{enumerate}
    \item Người dùng mở bộ thẻ ghi nhớ.
    \item Người dùng nhấn Xóa.
    \item Hệ thống hiển thị cảnh báo xác nhận (xóa sẽ xóa tất cả thẻ bên trong).
    \item Người dùng xác nhận xóa.
    \item Hệ thống xóa bộ thẻ và tất cả thẻ ghi nhớ liên quan.
    \item Người dùng được chuyển hướng về danh sách bộ thẻ.
\end{enumerate} \\
\hline
\textbf{Luồng thay thế} &
\begin{enumerate}
    \item Không có luồng thay thế.
\end{enumerate} \\
\hline
\textbf{Luồng ngoại lệ} &
\begin{enumerate}
    \item \textbf{Bộ thẻ không tồn tại:}
    \begin{enumerate}
        \item[1Err.] Hệ thống trả về lỗi 404.
    \end{enumerate}
    \item \textbf{Không có quyền:}
    \begin{enumerate}
        \item[1Err2.] Hệ thống trả về lỗi 403 khi người dùng không phải tác giả.
    \end{enumerate}
\end{enumerate} \\
\hline
\end{tabularx}
\captionof{table}{Đặc tả Use case: UC-0034 Xóa bộ thẻ ghi nhớ}
}
